\vspace{0.25cm}
Dans la partie « Rayonnement thermique », une étude qualitative du rayonnement du corps noir est proposée sans qu'aucune formule ne soit exigible. Celle-ci permet également d'aborder de manière quantitative l'effet de serre.
\begin{tabularx}{\textwidth}{|X|X|}
\hline 
\multicolumn{1}{|c|}{ Notions et contenus } & \multicolumn{1}{c|}{ Capacités exigibles } \\
\hline 
3.4. Rayonnement thermique & \\
\hline 
Approche descriptive du rayonnement du corps noir. & Exploiter les expressions fournies des lois de  Wien et de Stefan. \\
Loi de Wien, loi de Stefan. & Analyser quantitativement l'effet de serre en s'appuyant sur un bilan énergétique dans le cadre d'un modèle à une couche. \\
\hline
\end{tabularx}
